\chapter{TABELA DE SÍMBOLOS}

A Tabela de Símbolos é a estrutura de dados central encarregada de registrar e consultar identificadores de variáveis, funções, classes, métodos e atributos declarados ao longo da compilação de um programa em MiniPar 2026.2. 

No compilador, a gestão de símbolos está formalizada no componente \texttt{C\_TabelaSimbolos} (\texttt{src/componente\_tabela\_simbolos.py}), que realiza os contratos da interface \texttt{ITabelaSimbolos}.

\section{Estrutura de Dados e Pilha de Escopos Léxicos}

A arquitetura adota uma Pilha de Dicionários de Símbolos (\textit{Scope Stack}):
\begin{itemize}
    \item Escopo Global: corresponde à base da pilha (índice 0), contendo as declarações acessíveis universalmente em todo o módulo compilado.
    \item Escopos Locais Aninhados: a cada entrada em um novo contexto delimitado (corpo de funções, métodos de classes, blocos de controle \texttt{if}, laços \texttt{while}/\texttt{for} e blocos concorrentes \texttt{seq}/\texttt{par}), uma nova tabela local é empilhada através da operação \texttt{entrar\_escopo()}.
    \item Resolução de Identificadores (Sombra / \textit{Shadowing}): a consulta (\texttt{consultar(nome)}) percorre a pilha em sentido descendente (do topo atual para a base). Se um símbolo existir em um escopo mais interno, ele tem prioridade sobre definições em escopos externos, modelando o sombreamento de variáveis.
    \item Encerramento de Escopo: ao atingir o término de um bloco, o método \texttt{sair\_escopo()} desempilha a tabela local. Para alimentar relatórios didáticos na interface Web, o componente preserva um histórico de todos os símbolos registrados na sessão.
\end{itemize}

Cada símbolo é modelado como uma instância da classe \texttt{Symbol}, registrando:
\begin{itemize}
    \item \texttt{name}: Nome textual do identificador;
    \item \texttt{type}: Tipo associado (\texttt{int}, \texttt{number}, \texttt{string}, \texttt{bool}, \texttt{list}, \texttt{task}, nome da classe ou \texttt{var});
    \item \texttt{params}: Lista de tipos/nomes de parâmetros (no caso de funções e métodos);
    \item \texttt{attributes}: Dicionário de atributos de instância com seus respectivos tipos;
    \item \texttt{methods}: Dicionário contendo as assinaturas de métodos declarados na classe.
\end{itemize}

\section{Pseudocódigo da Tabela de Símbolos}

O Algoritmo \ref{algo:tabela_simbolos} formaliza a implementação do componente de Tabela de Símbolos, atendendo ao Requisito 6.

\begin{lstlisting}[language=Python, caption={Pseudocódigo da Tabela de Símbolos}, label={algo:tabela_simbolos}]
CLASSE Simbolo:
    ATRIBUTOS:
        nome: Texto
        tipo: Texto
        parametros: Lista de Textos <- []
        atributos: Dicionario[Texto, Texto] <- {}
        metodos: Dicionario[Texto, Lista de Textos] <- {}

CLASSE TabelaDeSimbolos:
    ATRIBUTOS:
        pilha_escopos: Pilha de Dicionarios[Texto, Simbolo]
        historico_declaracoes: Lista de Metadados
        proximo_id_escopo: Inteiro <- 1

    METODO inicializar():
        pilha_escopos <- Nova_Pilha()
        escopo_global <- Novo_Dicionario()
        pilha_escopos.empilhar(escopo_global)

    METODO entrar_escopo():
        novo_escopo_local <- Novo_Dicionario()
        pilha_escopos.empilhar(novo_escopo_local)
        proximo_id_escopo <- proximo_id_escopo + 1

    METODO sair_escopo():
        SE tamanho(pilha_escopos) > 1:
            pilha_escopos.desempilhar()

    METODO declarar(nome, tipo, linha=NULO, params=[], attrs={}, methods={}):
        escopo_atual <- pilha_escopos.topo()
        
        # Impede redeclaracao de identificador no mesmo escopo
        SE nome EXISTE_EM escopo_atual:
            RETORNAR FALSO  # Redeclaracao invalida detectada
        
        novo_simbolo <- Novo Simbolo(nome=nome, tipo=tipo, parametros=params, 
                                     atributos=attrs, metodos=methods)
        escopo_atual[nome] <- novo_simbolo

        # Registra no historico cumulativo para exibicao didatica
        registro <- Novo_Registro(nome=nome, tipo=tipo, linha=linha, 
                                  nivel_escopo=tamanho(pilha_escopos) - 1)
        historico_declaracoes.adicionar(registro)
        RETORNAR VERDADEIRO

    METODO consultar(nome):
        # Percorre a pilha do escopo mais interno (topo) para o global (base)
        PARA i DE (tamanho(pilha_escopos) - 1) ATE 0 PASSO -1:
            escopo <- pilha_escopos[i]
            SE nome EXISTE_EM escopo:
                RETORNAR escopo[nome]
        RETORNAR NULO  # Simbolo nao encontrado em nenhum escopo ativo

    METODO localizar_escopo(nome):
        PARA i DE (tamanho(pilha_escopos) - 1) ATE 0 PASSO -1:
            SE nome EXISTE_EM pilha_escopos[i]:
                RETORNAR i  # 0 indica global, >= 1 indica local
        RETORNAR NULO

    METODO listar_simbolos_historico():
        RETORNAR historico_declaracoes
FIM_CLASSE
\end{lstlisting}
